\chapter{Rational Maps}
\label{chp:rational_maps}

We present here some of the elementary properties of rational maps. A point
$\infty$ is adjoined to the complex plane and a rational map of degree $d$ is
considered as a $d$-fold map of the extended complex plane onto itself. The
spherical and chordal metrics are introduced, the important notions of
conjugacy, valency, fixed points and critical points are discussed, and the
Riemann-Hurwitz relation is proved.

% --- SECTION ------------------------------------------------------------------
% --- THE EXTENDED COMPLEX PLANE
% -----------------------------------------------
\section{The Extended Complex Plane}
\label{sec:the_extended_complex_plane}

\begin{definition}[Extended Complex Plane]
    The \emph{extended complex} plane is the union
    \begin{equation*}
        \ExtC = \C \cup \{\infty\},
    \end{equation*}
    where $\infty$ is an \enquote{abstract} point outside $\C$.
\end{definition}

Now let's define a metric in $\ExtC$ that handles all points with equal ease and
makes $\infty$ lose any special significance: \emph{the chordal metric}.

Let's start by defining the \emph{Riemann Sphere}. For that, we identify $\C$
with the horizontal plane
\begin{equation*}
    \{ (x_1, x_2, x_3) \in \R^3 \mid x_3 = 0 \}
\end{equation*}
in $\R^3$. Let $S$ be the unit sphere in $\R^3$ centered at the origin, and $N =
(0,0,1)$ be the north pole.

We now project each point in $\C$ to a point in $S$ by taking the intersection
of the unit sphere $S$ with the line passing through $N$ and the point $(x_1,
x_2, 0)$. An illustration of this projection is given in
Figure~\ref{fig:stereographic_projection}.

\begin{figure}[htbp]
    \centering
    \includegraphics[width=\textwidth]{beardons_book/stereographic_projection.png}
    \caption{Stereographic projection.}
    \label{fig:stereographic_projection}
\end{figure}

We identify this projection as:
\begin{equation*}
    \pi : \ExtC \longrightarrow S,
\end{equation*}
where for $z = x + \I y \in \C$:
\begin{equation*}
    \pi(z) = \left( \frac{2x}{1+|z|^2}, \frac{2y}{1+|z|^2}, \frac{|z|^2-1}{1+|z|^2} \right),
\end{equation*}
and $\pi(\infty) = (0,0,1)$.

\begin{definition}[Chordal Metric]
    Given the projection $\pi$ from $\ExtC$ to $S$, we define the \emph{chordal
    metric}, $\sigma : \ExtC \times \ExtC \rightarrow \R^+ \cup \{0\}$, as the
    Euclidean distance from $\pi(z)$ to $\pi(w)$:
    \begin{equation*}
        (z,w) \mapsto \sigma(z,w) = |\pi(z) - \pi(w)|.
    \end{equation*}
    Explicitly, when $z, w \in \C$:
    \begin{equation*}
        \sigma(z,w) = \frac{2|z-w|}{(1+|z|^2)^{1/2}(1+|w|^2)^{1/2}}.
    \end{equation*}
    For the point at infinity, we have:
    \begin{equation*}
        \sigma(z,\infty) = \lim_{w \rightarrow \infty} \sigma(z,w) = \frac{2}{(1+|z|^2)^{1/2}}.
    \end{equation*}
\end{definition}

\begin{remark}[$z \mapsto 1/z$ is an isometry]
    $\sigma(z,w) = \sigma(z^{-1}, w^{-1})$, i.e., the map $z \mapsto
    \frac{1}{z}$ is an isometry.
\end{remark}

\begin{remark}[Rotation in $S$ $\mapsto$ isometry in $\ExtC$]
    Any rotation $\varphi: S \rightarrow S$ of the sphere induces a map
    $\pi^{-1} \varphi \pi : \ExtC \rightarrow \ExtC$ which is an isometry. These
    isometries are Möbius transformations of the form:
    \begin{equation*}
        z \mapsto \frac{a z - \bar{c}}{c z + \bar{a}},
    \end{equation*}
    where $|a|^2 + |c|^2 = 1$.
\end{remark}

The chordal metric handles the point at infinity with ease, but it measures
distance through the three-dimensional space containing the Riemann sphere. If
we restrict our measurements strictly to the two-dimensional curved surface of
the sphere, we obtain the spherical metric.

\begin{definition}[Spherical Metric]
    Given the stereographic projection $\pi: \ExtC \to S$, the \emph{spherical
    distance} $\sigma_0(z,w)$ between two points in $\ExtC$ is defined as the
    Euclidean length of the shortest path (an arc of a great circle) connecting
    their projections $\pi(z)$ and $\pi(w)$ on the surface of $S$.
\end{definition}

If the chord joining the projected points subtends an angle $\theta$ at the
origin of the sphere, then by definition the spherical distance is exactly
$\sigma_0(z,w) = \theta$. By using the trigonometric connection to the chordal
metric, we can derive explicit formulas directly in terms of the complex
coordinates for $z, w \in \C$:
\begin{equation*}
    \sigma(z,w) = 2 \sin\left( \frac{\theta}{2} \right) = 2 \sin\left( \frac{\sigma_0(z,w)}{2} \right).
\end{equation*}
Thus,
\begin{align*}
    \sigma_0(z,w) &= 2 \arcsin\left( \frac{\sigma(z,w)}{2} \right) \\
                  &= 2 \arcsin\left( \frac{|z-w|}{(1+|z|^2)^{1/2}(1+|w|^2)^{1/2}} \right).
\end{align*}
For the point at infinity, taking the limit as $w \to \infty$ yields:
\begin{equation*}
    \sigma_0(z,\infty) = 2 \arcsin\left( \frac{1}{(1+|z|^2)^{1/2}} \right).
\end{equation*}

\begin{remark}[$\sigma \sim \sigma_0$]
    The chordal and spherical metrics are topologically equivalent ($\sigma(z,w)
    = 2 \sin(\theta/2)$). Furthermore, they satisfy the following inequalities:
    \begin{equation*}
        \frac{2}{\pi} \sigma_0(z,w) \le \sigma(z,w) \le \sigma_0(z,w)
    \end{equation*}
\end{remark}
\begin{proof}
    The inequalities follow directly from applying the elementary real-variable
    inequality $\frac{2\theta}{\pi} \le \sin\theta \le \theta$ on the interval
    $0 \le \theta \le \frac{\pi}{2}$ to the angle $\theta =
    \frac{\sigma_0(z,w)}{2}$.
\end{proof}

Because these metrics differ by at most a bounded scale factor of
$\frac{2}{\pi}$, we can use $\sigma$ and $\sigma_0$ interchangeably throughout
our study of rational dynamics whenever global constants are allowed to change.

\begin{remark}[Spherical length]
    Given a continuously differentiable curve $\gamma : I \rightarrow \ExtC$,
    its \emph{spherical length} is given by:
    \begin{equation*}
        \int_{\gamma} \frac{2|\D z|}{1+|z|^2}.
    \end{equation*}
\end{remark}

% --- SECTION ------------------------------------------------------------------
% --- RATIONAL MAPS
% ------------------------------------------------------------
\section{Rational Maps}
\label{sec:rational_maps}

\begin{definition}[Rational Map]
    A \emph{rational map} is a function of the form
    \begin{equation*}
        R(z) = \frac{P(z)}{Q(z)},
    \end{equation*}
    where $P(z)$ and $Q(z)$ are polynomials, not both being the zero polynomial.
\end{definition}

\begin{itemize}
    \item If $P = 0 \Rightarrow R = 0$.
    \item If $Q = 0 \Rightarrow R = \infty$.
    \item If $P(z_0) \neq 0$ and $Q(z_0) = 0$ for some $z_0 \in \C$, then
    $R(z_0) = \infty$.
    \item We define $R(\infty) = \lim_{z \rightarrow \infty} R(z)$.
    \item If $P(z_0) = 0$ and $Q(z_0) = 0$, strictly speaking $R(z_0)$ is not
    defined, but in this case we may cancel the common factors:
    \begin{align*}
        P(z) &= (z-z_0)^m P_1(z) \\
        Q(z) &= (z-z_0)^n Q_1(z)
    \end{align*}
    where $P_1(z_0) \neq 0$ and $Q_1(z_0) \neq 0$. Then:
    \begin{equation*}
        R(z) = \frac{P_1(z)}{Q_1(z)} (z-z_0)^{m-n}.
    \end{equation*}
\end{itemize}

From now on, we will consider that this has been done, i.e., $P$ and $Q$ are
coprime. Thus, given $R$, the polynomials $P$ and $Q$ are uniquely determined up
to scalar multiplication. This allows us to define the degree of a rational map.

\begin{definition}[Degree of a Rational Map, $\deg(R)$]
    Given a rational map $R = P/Q$ where $P$ and $Q$ are coprime polynomials,
    the \emph{degree of $R$} is defined as:
    \begin{equation*}
        \deg(R) = \max \{ \deg(P), \deg(Q) \}
    \end{equation*}
    where $\deg(\cdot)$ is the usual degree of a polynomial. We will consider
    that if $R = \alpha \in \ExtC$, then $\deg(R) = 0$.
\end{definition}

\begin{proposition}[{$\deg(RS) = \deg(R)\deg(S)$}]
    Given two rational function $R$ and $S$, the degree of $RS$ is $\deg(R)$
    times $\deg(S)$. That is,
    \begin{equation*}
        \deg(RS) = \deg(R)\deg(S).
    \end{equation*}
    Remember that by writing $RS$ we mean composition ($R \circ S$) not
    multiplication. However, $\deg(R)\deg(S)$ is a multiplication.
\end{proposition}
\begin{proof}
    See \cite[32]{Beardon1991}.
\end{proof}

\begin{corollary}[{$\deg(R^n) = [\deg(R)]^n$}] Given a rational function $R$,
    the degree of $R^n$ is the $n$-th power of $\deg(R)$. That is,
    \begin{equation*}
        \deg(R^n) = [\deg(R)]^n.
    \end{equation*}
\end{corollary}

Now, let's remember the definitions of \emph{holomorphic}, \emph{meromorphic}
and \emph{analytic} functions in $\C$ and see how they translate to $\ExtC$ with
the chordal metric.

\begin{definition}[Holomorphic]
    Given $D \subseteq \C$, a function $f : D \rightarrow \C$ is
    \emph{holomorphic} in $D$ if the derivative of $f$ exists at each point in
    $D$.
\end{definition}

\begin{definition}[Meromorphic]
    Given $D \subseteq \C$, a function $f : D \rightarrow \ExtC$ is
    \emph{meromorphic} in $D$ if each point in $D$ has a neighborhood on which
    either $f$ or $1/f$ is holomorphic.
\end{definition}

\begin{definition}[Pole of $f$]
    Given a meromorphic function $f$ in $D \subseteq \C$, the \emph{poles of
    $f$} are the points $w \in D$ where $f(w) = \infty$ ($\lim_{z \to w} f(z) =
    \infty$).
\end{definition}

\begin{remark}[$f$ meromorphic $\implies$ $1/f$ holomorphic at $f$'s poles]
    Given a meromorphic function $f$ in $D \subseteq \C$ with a pole at $w$,
    then $1/f$ is holomorphic near $w$ and $\frac{1}{f(w)} = 0$.
\end{remark}

\begin{definition}[$f$ is defined near $\infty$]
    A function $f$ is said to be \emph{defined near $\infty$} (or \emph{in a
    neighborhood of $\infty$}) if it is defined on some set $\{ z \in \C \mid
    |z| > r \} \cup \{\infty\}$.
\end{definition}

\begin{definition}[Holomorphic or Meromorphic at $\infty$]
    Given a function $f$ defined near $\infty$, $f$ is \emph{holomorphic} (or
    \emph{meromorphic}) \emph{at $\infty$} if $f(1/z)$ is holomorphic (or
    meromorphic) near $0$.
\end{definition}

\begin{definition}[Analytic in $D \subseteq \ExtC$]
    Given a map $f : D_1 \rightarrow D_2$ between any two subdomains of $\ExtC$,
    we say $f$ is \emph{analytic} if it is holomorphic or meromorphic at each
    point in $D_1$.
\end{definition}

\begin{remark}[$f$ meromorphic $\implies$ $f$ continuous at its poles]
    Given a meromorphic function $f$ in $D \subseteq \C$ with a pole at $w$,
    then (using the chordal metric and its topology), $f$ is continuous at the
    pole. This is because $z \mapsto 1/z$ is a $\sigma$-isometry:
    \begin{equation*}
        \sigma(f(z), f(w)) = \sigma\left(\frac{1}{f(z)}, \frac{1}{f(w)}\right) = \sigma\left(\frac{1}{f(z)}, 0\right) = \frac{2|1/f(z)|}{(1+|1/f(z)|^2)^{1/2}} \xrightarrow{z \rightarrow w} 0
    \end{equation*}
\end{remark}

Summing up, we have that the concepts of \emph{holomorphic}, \emph{meromorphic}
and \emph{analytic} functions translate well from $\C$ with the Euclidean
metric, to $\ExtC$ with the chordal or spheric metric. Hence, making $\infty$
lose any special significance it may have had in $\C$.

\begin{proposition}[Polynomials are analytic in $\ExtC$]
    Given a polynomial $P(z) = a_0 + a_1 z + \cdots + a_n z^n$ where $n > 0$ and
    $a_n \neq 0$. Then, $P$ is analytic in $\ExtC$.
\end{proposition}
\begin{proof}
    This proof serves to illustrate what holomorphic and meromorphic mean. We
    have to show that $P$ is holomorphic or meromorphic at each point in
    $\ExtC$. We have that $P$ is analytic in $\C$ because it's a polynomial,
    thus, we only have to show that $P$ is holomorphic or meromorphic at
    $\infty$. We'll see that it's meromorphic at $\infty$, i.e., $P(1/z)$ is
    meromorphic at $0$, i.e., $1/P(1/z)$ is holomorphic at $0$. And we can see
    that
    \begin{equation*}
        \frac{1}{P(1/z)} = \frac{z^n}{a_0 z^n + \cdots + a_n}
    \end{equation*}
    is differentiable at $0$, therefore, $P$ is meromorphic at $\infty$.
\end{proof}

In this proof we have also seen that all polynomials have a pole in $\infty$,
that is, $P(\infty) = \infty$. We will see that this means that $\infty$ is a
fixed point for all polynomials.

\begin{theorem}[Analytic in $\ExtC$ $\iff$ rational map]
    The rational maps are precisely the analytic maps of $\ExtC$ into itself.
\end{theorem}

\begin{definition}[$d$-fold map]
    We say a map $f: V \to W$ is a \emph{$d$-fold map} if, for every $w$ in $W$,
    the equation $f(z) = w$ has \emph{exactly} $d$ solutions in $V$ (counting
    multiple solutions by their multiplicity).
\end{definition}

\begin{theorem}[Rational maps are $d$-fold maps in $\ExtC$]
    \label{thm:rational_map_are_dfold_maps}
    Given $R$ a rational function of positive degree $d$. Then $R$ is a $d$-fold
    map of $\ExtC$ \emph{onto} itself.
\end{theorem}
\begin{proof}
    See \cite[31--32]{Beardon1991}.
\end{proof}

\begin{remark}[Zeros and poles of rational maps]
    Given $R=P/Q$ a rational function of positive degree $d$, where $P$ and $Q$
    are coprime polynomials with $\deg(P) = n$ and $\deg(Q) = m$. Then, the
    zeros of $R$ are the zeros of $P$, and the poles of $R$ are the zeros of
    $Q$. Moreover, the number of zeros and poles of $R$ at $\infty$ (counting
    multiplicities) are given by the following cases:
    \begin{enumerate}[label=(\roman*)]
        \item if $n = m$, then all the zeros and poles of $R$ are in $\C$;
        \item if $n > m$, then the number of poles at $\infty$ is $n-m$;
        \item and if $n < m$, then the number of zeros at $\infty$ is $m-n$.
    \end{enumerate}
\end{remark}

% --- SECTION ------------------------------------------------------------------
% --- THE LIPSCHITZ CONDITION
% --------------------------------------------------
\section{The Lipschitz Condition}
\label{sec:the_lipschitz_condition}

\begin{theorem}[Rational maps satisfy a Lipschitz condition]
    \label{thm:rational_maps_satify_a_lipschitz_condition}
    A rational map $R$ satisfies some \emph{Lipschitz condition}
    \begin{equation*}
        \sigma_0(R(z), R(w)) \leq M\sigma_0(z,w)
    \end{equation*}
    on the complex sphere.
\end{theorem}
\begin{proof}
    See \cite[33]{Beardon1991}.
\end{proof}

\begin{theorem}[Möbius transformations explicit Lipschitz condition]
    The Möbius map
    \begin{equation*}
        g(z) = \frac{az + b}{cz+d},\quad ad-bc = 1,
    \end{equation*}
    satisfies the \emph{Lipschitz condition}
    \begin{equation*}
        \sigma_0(g(z), g(w)) \leq ||g||^2 \sigma_0(z, w).
    \end{equation*}
\end{theorem}
\begin{proof}
    See \cite[33--34]{Beardon1991}.
\end{proof}

\begin{theorem}[Möbius transformations family explicit uniform Lipschitz condition]
    Let $m$ be any positive number. Then the Möbius transformations $g$ which
    satisfy
    \begin{alignat*}{3}
        &\sigma(g(0),      &&g(1))      &&\geq m, \\
        &\sigma(g(1),      &&g(\infty)) &&\geq m, \\
        &\sigma(g(\infty), &&g(0))      &&\geq m
    \end{alignat*}
    also satisfy the \emph{uniform Lipschitz condition}
    \begin{equation*}
        \sigma(g(z), g(w)) \leq \frac{\pi}{m^3} \sigma(z, w).
    \end{equation*}
\end{theorem}
\begin{proof}
    See \cite[34--35]{Beardon1991}.
\end{proof}

\begin{definition}[Complementary components of a closed curve $\gamma$]
    Let $\gamma$ be a \emph{closed curve} on the sphere. The complement of
    $\gamma$ is a union of mutually disjoint domains, which we call the
    \emph{complementary components} of $\gamma$.
\end{definition}

\begin{definition}[Exterior and interior of a closed curve $\gamma$ contained in a open hemisphere]
    Let $\gamma$ be a \emph{closed curve} that lies in some \emph{open
    hemisphere}. Then exactly one of the complementary components contains a
    hemisphere; we denote this by $E(\gamma)$ and call it the \emph{exterior} of
    $\gamma$. The other complementary components of $\gamma$ are said to be the
    \emph{interior components} of $\gamma$.
\end{definition}

\begin{theorem}[Interior components mapping theorem for rational maps]
    \label{thm:iterior_components_mapping_theorem_for_rational_maps}
    Let $R$ be a rational function. Then there is a positive number $\delta$
    such that if $\gamma$ is any closed curve of $\sigma_0$-diameter less than
    $\delta$, then the image $R(\Omega)$ of an interior component $\Omega$ of
    $\gamma$ does not meet the exterior $E(R(\gamma))$ of $R(\gamma)$.
\end{theorem}
\begin{proof}
    See \cite[35--36]{Beardon1991}.
\end{proof}

This result almost says that if the curve $\gamma$ is sufficiently small, then
the interior components of $\gamma$ map to the interior components of the image
curve $R(\gamma)$: however, the images of the interior components may also
contain points of $R(\gamma)$, so it is necessary to express this a little more
carefully in terms of the exterior of the image curve
(Figure~\ref{fig:interior_of_curve_mobius_transformed}).

\begin{figure}[htbp]
    \centering
    \resizebox{\textwidth}{!}{\begin{tikzpicture}[
    axis/.style={-{Stealth[length=2.5mm]}},
    thick line/.style={thick, line cap=round, line join=round}
]

    % --- LEFT PLOT: \Omega and \gamma -----------------------------------------
    
    % Shape Path Definition
    \draw[thick] (0,0)
        to[out=120, in=0] (-1,0.75)
        to[out=-180, in=90] (-2,0)
        to[out=-90, in=180] (-1,-0.75)
        to[out=0, in=-120] (0,0);

    \draw[thick, fill=gray!30] (0,0)
        to[out=60, in=180] (1,0.75)
        to[out=0, in=90] (2,0)
        to[out=-90, in=0] (1,-0.75)
        to[out=180, in=-60] (0,0);

    % Axes
    \draw[axis] (-2.5,0) -- (2.5,0);
    \draw[axis] (0,-2.5) -- (0,2.5);
    
    % Labels
    \node at (1, 0.25) {$\Omega$};
    \node at (-1.75, 0.75) {$\gamma$};

    % --- MIDDLE ARROW: R(z) ---------------------------------------------------
    \draw[-{Stealth[length=3mm]}, very thick] (3.0,0) -- (4.2,0) node[midway, above=2pt] {\large $R$};

    % --- RIGHT PLOT: R(\Omega) and R(\gamma)  ---------------------------------
    \begin{scope}[shift={(7.5,0)}]
        
        \begin{scope}[scale=0.45, transform shape, shift={(0, 0)}]
            \begin{pgfonlayer}{nodelayer}
                \node (0) at (1, 1) {};
                \node (1) at (1, 2) {};
                \node (2) at (1.5, 3) {};
                \node (3) at (0, 4) {};
                \node (4) at (-1, 3) {};
                \node (5) at (0, 2) {};
                \node (6) at (2.25, 1.75) {};
                \node (7) at (-1, 1.25) {};
                \node (8) at (-3, 0.5) {};
                \node (9) at (-2, -1) {};
                \node (10) at (-1.5, -2) {};
                \node (11) at (-2.5, -2.75) {};
                \node (12) at (-3.25, -1.75) {};
                \node (13) at (0, -1.75) {};
                \node (14) at (2, -2) {};
                \node (15) at (3.25, -1) {};
                \node (16) at (2.75, 0) {};
                \node (17) at (2.25, 1.25) {};
                \node (18) at (-3, -0.25) {};
            \end{pgfonlayer}

            % 1. Left Top Loop (Nodes: 1 -> 2 -> 3 -> 4 -> 5 -> 1)
            \fill[gray!30] (1.center) 
                [bend right] to (2.center)
                [bend right=45, looseness=0.75] to (3.center)
                [bend right] to (4.center)
                [bend right=60] to (5.center)
                -- cycle; % '-- cycle' cleanly closes the path back to (1.center)

            % 2. Right Top Loop (Nodes: 0 -> 1 -> 6 -> 17 -> 0)
            \fill[gray!30] (0.center) 
                [bend left=60] to (1.center)
                [bend left=15] to (6.center)
                [bend left=60, looseness=1.25] to (17.center)
                [bend left=15, looseness=0.75] to cycle;

            \begin{pgfonlayer}{edgelayer}
                \draw [bend left=60, thick] (0.center) to (1.center);
                \draw [bend right, thick] (1.center) to (2.center);
                \draw [bend right=45, looseness=0.75, thick] (2.center) to (3.center);
                \draw [bend right, thick] (3.center) to (4.center);
                \draw [bend right=45, thick] (4.center) to (5.center);
                \draw [thick] (5.center) to (1.center);
                \draw [bend left=15, thick] (1.center) to (6.center);
                \draw [thick] (0.center) to (7.center);
                \draw [bend left=330, looseness=0.75, thick] (7.center) to (8.center);
                \draw [bend left, thick] (9.center) to (10.center);
                \draw [bend left=45, thick] (10.center) to (11.center);
                \draw [bend left=45, thick] (11.center) to (12.center);
                \draw [bend left=45, thick] (12.center) to (9.center);
                \draw [bend left=15, thick] (9.center) to (13.center);
                \draw [bend right, thick] (13.center) to (14.center);
                \draw [bend right=15, thick] (14.center) to (15.center);
                \draw [bend right=45, thick] (15.center) to (16.center);
                \draw [bend left=15, looseness=0.75, thick] (16.center) to (0.center);
                \draw [bend left=60, looseness=1.25, thick] (6.center) to (17.center);
                \draw [bend left=15, looseness=0.75, thick] (17.center) to (0.center);
                \draw [bend right, thick] (8.center) to (18.center);
                \draw [bend right=15, looseness=0.75, thick] (18.center) to (9.center);
            \end{pgfonlayer}
        \end{scope}
     
        % Axes
        \draw[axis] (-2.5,0) -- (2.5,0);
        \draw[axis] (0,-2.5) -- (0,2.5);

        % Labels
        \node at (1.25, 1.25) {$R(\Omega)$};
        \node at (-1.5, 0.7) {$R(\gamma)$};
        
    \end{scope}

\end{tikzpicture}}
    \caption{If the curve $\gamma$ is sufficiently small
    ($\diameter_{\sigma_0}(\gamma) < \delta$), the image of its interior
    component $\Omega$ under a rational function $R$ is entirely contained
    within the bounded regions defined by the image curve $R(\gamma)$, never
    entering the exterior $E(R(\gamma))$.}
    \label{fig:interior_of_curve_mobius_transformed}
\end{figure}

% --- SECTION ------------------------------------------------------------------
% --- CONJUGACY
% ----------------------------------------------------------------
\section{Conjugacy}
\label{sec:conjugacy}

\begin{definition}[Conjugate rational maps]
    Given two rational maps $R$ and $S$, we say $R$ and $S$ are \emph{conjugate}
    if there is some Möbius map g with
    \begin{equation*}
        S = g R g^{-1}.
    \end{equation*}
\end{definition}

\begin{remark}[Conjugacy is an equivalence relation]
    Conjugacy is an equivalence relation, and the equivalence classes are
    \emph{conjugacy classes} of rational maps.
\end{remark}

\begin{remark}[Invariant properties of conjugacy]
    Given two conjugate rational maps $R$ and $S$, where $S = g R g^{-1}$. Then,
    the following are true:
    \begin{enumerate}[label=(\roman*)]
        \item $\deg(R) = \deg(S)$;
        \item $S^n = g R^n g^{-1}$ (this means that we can transfer a problem
        concerning $R$ to a (possibly simpler) problem concerning a conjugate
        $S$);
        \item and $S$ fixes $g(z)$ if and only if $R$ fixes $z$ (conjugacy
        respects fixed points).
    \end{enumerate}
\end{remark}

All of these observations are elementary; nevertheless, they are important
because, in general, we shall not bother to distinguish between conjugate
rational functions.

As a simple illustration of the use of conjugacy, let us characterize the
polynomials within the class of rational maps.

\begin{remark}[Characterization of polynomials]
    A rational map $R$ is a polynomial if and only if $R$ has a pole at $\infty$
    and no poles in $\C$; more succinctly, if and only if $R^{-1}\{\infty\} =
    \{\infty\}$.
\end{remark}

\begin{theorem}[Characterization of the conjugacy class of the polynomials]
    \label{thm:characterization_of_the_conjugacy_class_of_the_polynomials}
    A non-constant rational map $R$ is conjugate to a polynomial if and only if
    there is some $w$ in $\ExtC$ with $R^{-1}\{w\} = \{w\}$.
\end{theorem}

% --- SECTION ------------------------------------------------------------------
% --- VALENCY
% ------------------------------------------------------------------
\section{Valency}
\label{sec:valency}

Given any function $f$ that is non-constant and holomorphic near the point $z_0$
in $\C$. Then $f$ has a Taylor expansion at $z_0$, say
\begin{equation*}
    f(z) = f(z_0) + a_k(z-z_0)^k + a_{k+1}(z-z_0)^{k+1} + \dots
\end{equation*}
where $a_k \neq 0$ and the positive integer $k$ is uniquely determined by the
condition that the limit
\begin{equation*}
    \lim_{z \to z_0} \frac{f(z)-f(z_0)}{(z-z_0)^k}
\end{equation*}
exists, is finite, and is non-zero. 

\begin{definition}[Valency of $f$ at $z_0 \in \C$]
    \label{def:valency_in_C}
    We denote this integer $k$ by $v_f(z_0)$ and call it the \emph{valency}, or
    \emph{order}, of $f$ at $z_0$.
\end{definition}

\begin{remark}[Valency is a local property]
    Note that when defining the valency, we've used the Taylor expansion at
    $z_0$. This expansion is defined \emph{near the point $z_0$}, that is,
    valency is defined locally in a neighborhood of a point.
\end{remark}

\begin{remark}[$v_f(z_0)$ is the number of solutions of $f(z) = f(z_0)$]
    Of course, the valency of $f$ is the number of solutions of $f(z) = f(z_0)$
    at $z_0$ (restricting $f$ to a neighborhood near $z_0$).
\end{remark}

\begin{proposition}[Chain Rule for Valency]
    The valency function $v$ satisfies the important \emph{Chain Rule}:
    \begin{equation*}
        v_{fg}(z_0) = v_f(g(z_0))v_g(z_0),
    \end{equation*}
    where $f$ and $g$ are non-constant holomorphic functions near $g(z_0)$ and
    $z_0$ respectively; and $z_0$, $g(z_0)$, and $fg(z_0)$ are all in $\C$.
\end{proposition}
\begin{proof}
    See \cite[37--38]{Beardon1991}.
\end{proof}

\begin{remark}[Local injectivity in terms of valency]
    \label{rmk:local_injectivity_in_terms_of_valency}
    Given any function $f$ that is non-constant and holomorphic near the point
    $z_0$ in $\C$. Observe that $f$ is injective in some neighborhood of $z_0$
    if and only if $v_f(z_0) = 1$.
\end{remark}

\begin{remark}[Composing with injective functions preserves valency]
    Because of the previous remark
    (\ref{rmk:local_injectivity_in_terms_of_valency}), the Chain Rule shows that
    the valency is preserved under a pre-application, and a post-application, of
    an injective function.

    Explicitly, if $fg$ is defined near $z$ with $g$ injective near $z$, and if
    $hf$ is defined near $z$ with $h$ injective near $f(z)$, then:
    \begin{align*}
        v_{fg}(z) &= v_f(g(z)), \\
        v_{hf}(z) &= v_f(z).
    \end{align*}
\end{remark}

The properties of valency under injective functions enable us to extend our
definition of $v_f(z_0)$ to the case where $z_0 = \infty$, or $f(z_0) = \infty$
(or both).

\begin{definition}[Valency at $\infty$]
    \label{def:valency_in_ExtC}
    Given any function $f$ that is non-constant and holomorphic near the point
    $\infty$ in $\ExtC$ or $z_0$ in $\ExtC$ such that $f(z_0) = \infty$ (or
    both). Then, we can select any Möbius maps $g$ and $h$ which map $z_0$ and
    $f(z_0)$ respectively to points in $\C$ and then define $v_f(z_0)$ by
    \begin{equation*}
        v_f(z_0) = v_F(g(z_0)),
    \end{equation*}
    where $F = hfg^{-1}$.
\end{definition}

The essential observation is that if we repeat this procedure using different
maps $g$ and $h$, then the resulting answer will be the same; thus $v_f(z_0)$ is
defined independently of the choice of $g$ and $h$.

\begin{definition}[Univalent map]
    An analytic map $f : D \to \ExtC$ is said to be \emph{univalent} in $D$ if
    it is injective there.
\end{definition}

\begin{remark}[$f$ univalent $\implies v_f(z) = 1 \forall z \in D$]
    If $f$ is univalent in $D$, then $v_f(z) = 1$ for every $z$ in $D$.
\end{remark}

The converse \emph{need not hold}, being univalent is a global property (the
whole domain of definition of the function), but valency is a local property (a
neighborhood of $z_0$).

\begin{example}
    The map $z \mapsto z^2$ of $\C - \{0\}$ onto itself has valency $1$
    everywhere, but it is not injective in $\C - \{0\}$.
\end{example}

\begin{remark}[Rational maps are $d$-fold maps via valency]
    We've seen that rational maps are $d$-fold maps
    (Theorem~\ref{thm:rational_map_are_dfold_maps}), now we can write this in
    terms of valency. Given a non-constant rational map $R$ of degree $d$, then
    for any $z_0$ in $R^{-1}\{w\}$, $v_R(z_0)$ is the number of solutions of
    $R(z) = w$ at $z_0$, and so for each $w$ in the complex sphere,
    \begin{equation*}
        \sum_{z \in R^{-1}\{w\}} v_R(z) = \deg(R).
    \end{equation*}
\end{remark}

% --- SECTION ------------------------------------------------------------------
% --- THE INVERSE FUNCTION THEOREM
% ---------------------------------------------
\section{The Inverse Function Theorem}
\label{sec:the_inverse_function_theorem}

Now that we have established the local behavior of analytic maps via valency, we
can formalize when these maps can be locally inverted. In the classical complex
plane, this is governed by the derivative.

\begin{theorem}[Inverse Function Theorem in $\C$]
    \label{thm:inverse_function_theorem_in_C}
    Let $f: U \to \C$ be a holomorphic function on an open set $U \subseteq \C$,
    and let $z_0 \in U$. If $f'(z_0) \neq 0$, then there exist open
    neighborhoods $V \subseteq U$ of $z_0$ and $W$ of $f(z_0)$ such that the
    restriction $f: V \to W$ is a bijection.
    
    Furthermore, the local inverse function $f^{-1}: W \to V$ is holomorphic,
    and its derivative satisfies:
    \begin{equation*}
        (f^{-1})'(f(z_0)) = \frac{1}{f'(z_0)}.
    \end{equation*}
\end{theorem}

Let us now bridge this classical theorem with our structural definition of
valency at a point.

\begin{remark}[Local injectivity in terms of derivative and valency in $\C$]
    By examining the Taylor expansion of a non-constant holomorphic function $f$
    near $z_0 \in \C$:
    \begin{equation*}
        f(z) = f(z_0) + f'(z_0)(z-z_0) + \frac{f''(z_0)}{2!}(z-z_0)^2 + \cdots
    \end{equation*}
    we see that $f'(z_0) \neq 0$ if and only if the first non-zero coefficient
    appears at index $k=1$. By Definition~\ref{def:valency_in_C}, this means
    $v_f(z_0) = 1$.
    
    Thus, the Inverse Function Theorem in $\C$
    (Theorem~\ref{thm:inverse_function_theorem_in_C}) states that $f$ possesses
    a local holomorphic inverse near $z_0 \in \C$ if and only if $v_f(z_0) = 1$.
\end{remark}

Just as we did with valency, we can exploit the fact that composing with
injective functions preserves local bijectivity to extend this theorem to cases
where $z_0 = \infty$, $f(z_0) = \infty$, or both.

\begin{theorem}[Inverse Function Theorem in $\ExtC$]
    \label{thm:inverse_function_theorem_in_ExtC}
    Let $f : D_1 \to D_2$ be a non-constant analytic map between two subdomains
    of $\ExtC$, and let $z_0 \in D_1$. Then, we can select any two Möbius
    transformations $g$ and $h$ which map $z_0$ and $f(z_0)$ respectively to
    points in $\C$. Let $F = hfg^{-1}$ be the localized map defined near $g(z_0)
    \in \C$. If $F'(g(z_0)) \neq 0$, then $f$ is a local bijection from a
    $\sigma$-neighborhood of $z_0$ onto a $\sigma$-neighborhood of $f(z_0)$.
    Consequently, it possesses a local analytic inverse $f^{-1}$ near $f(z_0)$.
\end{theorem}

The essential observation is that if we repeat this procedure using different
maps $g$ and $h$, then the resulting answer will be the same; thus the theorem
can be generalized to $\ExtC$ from $\C$ independently of the choice of $g$ and
$h$.

\begin{remark}[Local injectivity in terms of localized derivative and valency in $\ExtC$]
    Let $f$, $g$, $h$, and $F$ be defined as in
    Theorem~\ref{thm:inverse_function_theorem_in_ExtC}. By expanding the
    non-constant holomorphic map $F$ into a Taylor series near the point $w_0 =
    g(z_0) \in \C$:
    \begin{equation*}
        F(w) = F(w_0) + F'(w_0)(w-w_0) + \frac{F''(w_0)}{2!}(w-w_0)^2 + \cdots
    \end{equation*}
    we observe that $F'(g(z_0)) \neq 0$ if and only if the first non-zero
    coefficient after the constant term appears at index $k=1$. 
    
    By your definition of extended valency, this means $v_F(g(z_0)) = v_f(z_0) =
    1$. Thus, the Inverse Function Theorem in $\ExtC$
    (Theorem~\ref{thm:inverse_function_theorem_in_ExtC}) establishes a universal
    rule across the entire Riemann sphere: an analytic map $f$ admits a local
    analytic inverse near a point $z_0 \in \ExtC$ if and only if $v_f(z_0) = 1$.
\end{remark}

% --- SECTION ------------------------------------------------------------------
% --- FIXED POINTS
% -------------------------------------------------------------
\section{Fixed Points}
\label{sec:fixed_points}

\begin{definition}[Fixed point of $f$]
    Given any function $f : D \to \ExtC$ defined on some domain $D \subseteq
    \ExtC$. A point $\zeta$ in $D$ is a \emph{fixed point} of $f$ if $f(\zeta) =
    \zeta$.
\end{definition}

The fixed points of a mapping play an important role in much of our later work,
and they will be studied and classified in Chapter~\ref{chp:periodic_points}.

\begin{proposition}[Fixed points of rational maps]
    Given by two coprime polynomials $P$ and $Q$ and the rational map $R = P/Q$.
    Then,
    \begin{enumerate}[label=(\roman*)]
        \item $R$ fixes $\infty$ if and only if $\deg(P) > \deg(Q)$;
        \item and $R$ fixes $\zeta \in \C$ if and only if $P(\zeta) = \zeta
        Q(\zeta)$; that is, the \emph{fixed points of $R$ in $\C$} are the
        solutions of 
        \begin{equation*}
            P(z) - z Q(z) = 0.
        \end{equation*}
    \end{enumerate}
\end{proposition}
\begin{proof}
    See \cite[39]{Beardon1991}.
\end{proof}

We count fixed points the same way we count zeros, that is, counting
multiplicities.

\begin{definition}[Number of fixed points at $\zeta \in \C$]
    Given any function $f : D \to \ExtC$ defined on some domain $D \subseteq
    \ExtC$. The number of fixed points of $f$ at a point $\zeta$ in $\C$ is
    given by the valency of the function $f(z) - z$ at $\zeta$, i.e.,
    $v_{f(z)-z}(\zeta)$.
\end{definition}

\begin{example}
    The map $f(z) = z + z^3$ has three fixed points at $0$. This is because the
    fixed points of $f$ are the zeros of $f(z) - z$, and
    \begin{equation*}
        f(z) - z = z^3,
    \end{equation*}
    which has valency $3$ at $0$ ($z^3 = f(0) - 0 = 0$ has three solutions).
    This is not to be confused with the valency of $f$ at $0$, which is $1$. In
    fact, $f(z) = f(0)$ has only one solution at $0$ ($f(z) = z(1 + z^2)$).
\end{example}

To count the number of fixed points at $\infty$, we need

\begin{lemma}[$\zeta \in \C$ fixed by $f \implies \zeta$ fixed by \enquote{conjugate} in $\C$]
    \label{lem:number_of_fixed_points_is_preserved_under_conjugacy_in_C}
    Let $\zeta$ in $\C$ be a fixed point of an analytic map $f$, and let
    $\varphi$ be any map that is analytic, injective and finite in some
    neighborhood of $\zeta$. Then, $\varphi f \varphi^{-1}$ has the same number
    of fixed points at $\varphi(\zeta)$ as $f$ has at $\zeta$.
\end{lemma}
\begin{proof}
    See \cite[39--40]{Beardon1991}.
\end{proof}

\begin{definition}[Number of fixed points at $\infty$]
    Given any function $f : D \to \ExtC$ defined on some domain $D \subseteq
    \ExtC$, $\infty \in D$. Then we can select an analytic map $\varphi$, that
    is injective and finite in some neighborhood of $\infty$, and then define
    the number of fixed points of $f$ at $\infty$ by the number of fixed points
    of $\varphi f \varphi^{-1}$ at $\varphi(\infty)$.
\end{definition}

The essential observation is that if we repeat this procedure using different
maps $\varphi$, then the resulting answer will be the same; thus, the number of
fixed points of $f$ at $\infty$ is defined independently of the choice of
$\varphi$.

With this definition,
Lemma~\ref{lem:number_of_fixed_points_is_preserved_under_conjugacy_in_C} implies

\begin{theorem}[Number of fixed points is preserved under conjugacy]
    Given two conjugate rational maps $R$ and $S$, where $S = g R g^{-1}$, $g$
    is a Möbius transformation. Then, the number of fixed points of $R$ at
    $\zeta$ is the same as the number of fixed points of $S$ at $g(\zeta)$.
\end{theorem}

\begin{theorem}[Number of fixed points of a rational map]
    Given a rational map $R$ of degree $d \geq 1$. Then, the number of fixed
    points of $R$ in $\ExtC$ is exactly $d + 1$ (counting multiplicities).
\end{theorem}
\begin{proof}
    See \cite[40]{Beardon1991}.
\end{proof}

\begin{definition}[Multiplier of a fixed point in $\C$]
    Given a rational map $R$ and a fixed point $\zeta$ in $\C$ of $R$. We define
    the \emph{multiplier} of $R$ at $\zeta$, denoted by $m(R, \zeta)$, as the
    complex number given by the derivative:
    \begin{equation*}
        m(R, \zeta) = R'(\zeta).
    \end{equation*}
\end{definition}

\begin{remark}[Multiplier in $\C$ is preserved under conjugation]
    Note that since the derivative is invariant under a conjugation, the
    multiplier is also invariant under a conjugation (provided the corresponding
    point is also in $\ExtC$).
\end{remark}

\begin{definition}[Multiplier of the fixed point $\infty$]
    Given a rational map $R$ with a fixed point at $\infty$. Then, we can select
    a Möbius map $g$ that maps $\infty$ to some point in $\C$, and then define
    the multiplier of $R$ at $\infty$ by the multiplier of $g R g^{-1}$ at
    $g(\infty)$:
    \begin{equation*}
        m(R, \infty) = m(g R g^{-1}, g(\infty)).
    \end{equation*}
\end{definition}

Note that, by the preceding remarks, this definition is independent of the
choice of $g$.

\begin{remark}[Multiplier is preserved under conjugation]
    We now have defined a \emph{conjugation invariant multiplier $m(R, \zeta)$}
    of any rational map $R$ at any of its fixed points.
\end{remark}

\begin{example}[Multiplier at $\infty$ for a rational map]
    Let $R$ be a rational map of the form
    \begin{equation*}
        R(z) = \frac{a_0 + a_1 z + \cdots + a_n z^n}{b_0 + b_1 z + \cdots + b_m z^m},
    \end{equation*}
    where $a_n b_m \neq 0$ and $n > m$, so $R$ fixes $\infty$. By definition,
    $m(R, \infty)$ is the derivative of the conjugate map $S(z) = 1/R(1/z)$ at
    the origin. A direct computation shows:
    \begin{equation*}
        S'(0) =
        \begin{cases}
            b_m/a_n & \text{if } n = m + 1, \\
            0 & \text{if } n > m + 1.
        \end{cases}
    \end{equation*}
    On the other hand, computing the derivative $R'(z)$ and taking the limit
    gives:
    \begin{equation*}
        R'(\infty) = \lim_{z \to \infty} R'(z) =
        \begin{cases}
            a_n/b_m & \text{if } n = m + 1, \\
            \infty & \text{if } n > m + 1.
        \end{cases}
    \end{equation*}
\end{example}

\begin{remark}[Formula for the multiplier of a rational map at a fixed point in $\ExtC$]
    These results show that the multiplier $m(R, \zeta)$ of a rational map $R$
    at fixed point $\zeta$ is given by
    \begin{equation*}
        m(R, \zeta) =
        \begin{cases}
            R'(\zeta) & \text{if } \zeta \neq \infty; \\
            \frac{1}{R'(\infty)} & \text{if } \zeta = \infty.
        \end{cases}
    \end{equation*}
\end{remark}

\begin{example}
    For example, if $R(z) = 3z$, then $\infty$ is an attracting fixed point of
    $R$ (this is geometrically clear) and in this case, $m(R, \infty) = 1/3$.
\end{example}

We conclude this section by analyzing the behavior of the number of fixed points
under iteration near a fixed point of $f$.

\begin{theorem}[Iterates near the fixed point $0$]
    Given an analytic function $f$ that in a neighborhood of the origin has the
    expansion
    \begin{equation*}
        f(z) = a z + b_1 z^{r+1} + \cdots,
    \end{equation*}
    where $a \neq 0$, $b_1 \neq 0$, and $r \geq 1$. Then its $n$-th iterate, in
    that neighborhood, has the form
    \begin{equation*}
        f^n(z) = a^n z + b_n z^{r+1} + \cdots,
    \end{equation*}
    where
    \begin{equation*}
        b_n = a^{n-1} b_1 \left( 1 + a^r + a^{2r} + \cdots + a^{(n-1)r} \right).
    \end{equation*}
\end{theorem}
\begin{proof}
    See \cite[42]{Beardon1991}.
\end{proof}

From this theorem, we immediately obtain several important corollaries regarding
the behavior of coefficients and the multiplicity of fixed points under
iteration.

\begin{corollary}[$a = 1 \implies b_n = n b_1$]
    \label{cor:a_equals_1_implies_bn_equals_nb1}
    If $a = 1$, then $b_n = n b_1$.
\end{corollary}
\begin{proof}
    This is a direct consequence of the formula for $b_n$:
    \begin{equation*}
        b_n = a^{n-1} b_1 \left( 1 + a^r + a^{2r} + \cdots + a^{(n-1)r} \right) = b_1 (1 + 1 + \cdots + 1) = n b_1.
    \end{equation*}
\end{proof}

\begin{corollary}[$b_n = 0 \iff a^r \neq 1$ and $a^{nr} = 1$]
    $b_n = 0$ if and only if $a^r \neq 1$ and $a^{nr} = 1$.
\end{corollary}
\begin{proof}
    The terms inside the parentheses in the formula for $b_n$ form a geometric
    series. Thus, we can write:
    \begin{equation*}
        b_n = a^{n-1} b_1 \frac{1 - a^{nr}}{1 - a^r}.
    \end{equation*}
    For $b_n$ to be zero, the numerator of that fraction must be zero, but the
    denominator must not be zero (to avoid a $0/0$ situation, which we already
    showed in the previous corollary
    (Corollary~\ref{cor:a_equals_1_implies_bn_equals_nb1}) that it evaluates to
    $n$).
    \begin{alignat*}{4}
        &1 - a^{nr} &&= 0    &&\implies a^{nr} &&= 1, \\
        &1 - a^r    &&\neq 0 &&\implies a^r    &&\neq 1.
    \end{alignat*}
    Thus, $b_n = 0$ if and only if $a^r \neq 1$ and $a^{nr} = 1$.
\end{proof}

\begin{remark}[$f^n$ has a single fixed point at the origin $\implies f$ has a single fixed point at the origin]
    If $a^n \neq 1$, then $a \neq 1$. Thus, if $f^n$ has a single fixed point at
    the origin, then so does $f$.
\end{remark}
\begin{proof}
    If $a = 1$, then $a^n = 1$ for all $n$. Thus, if $a^n \neq 1$, then $a \neq
    1$.

    If $a \neq 1$, the function $f(z) - z$ looks like $(a-1)z + b_1 z^{r+1} +
    \cdots = z((a-1) + b_1 z^r + \cdots)$, where $a-1 \neq 0$. Thus, $f(z)$ has
    a single fixed point at the origin.

    If $f^n$ has a single fixed point at the origin, then we can only factor out
    one $z$ from $f^n(z) - z$, giving us $f^n(z) - z = (a^n - 1)z + b_n z^{r+1}
    + \cdots = z((a^n - 1) + b_n z^r + \cdots)$, where $a^n - 1 \neq 0$.

    Combining these three results, we have that if $f^n$ has a single fixed
    point at the origin, then $a^n \neq 1$, which implies $a \neq 1$, thus $f$
    has a single fixed point at the origin.
\end{proof}

\begin{corollary}[Number of fixed points at $0$ of the iterates]
    \label{cor:fn_has_at_least_as_many_fixed_points_as_f}
    The iterate $f^n$ has at least as many fixed points at the origin as $f$
    has. Furthermore, if $f^n$ has strictly more fixed points at the origin than
    $f$, then $a \neq 1$ but $a^n = 1$.
\end{corollary}
\begin{proof}
    Let's begin by showing that $f^n$ has at least as many fixed points at the
    origin as $f$ has. We have two cases:
    \begin{itemize}
        \item If $a \neq 1$: $f(z) - z$ has a single fixed point at the origin.
        Since $f^n(0) = 0$, $0$ must be a fixed point of $f^n$ of multiplicity
        at least $1$.
        \item If $a = 1$: then $f(z) - z = b_1 z^{r+1} + \cdots$, so $0$ is a
        fixed point of $f$ of multiplicity $r + 1$ ($b_1 \neq 0$). On the other
        hand, by Corollary~\ref{cor:a_equals_1_implies_bn_equals_nb1}, $f^n(z) -
        z = n b_1 z^{r+1} + \cdots$, so $0$ is a fixed point of $f^n$ of
        multiplicity $r + 1$ (the same as for $f$). Thus, if $a = 1$, then $f^n$
        has as many fixed points at the origin as $f$ has.
    \end{itemize}

    Now, let's show that if $f^n$ has strictly more fixed points at the origin
    than $f$, then $a \neq 1$ but $a^n = 1$. From our analysis above, if $a =
    1$, the multiplicities are exactly equal, so for $f^n$ to have strictly more
    fixed points at the origin than $f$, we cannot have $a = 1$. Thus, $a \neq
    1$.

    If $a \neq 1$, then $f$ has a single fixed point at the origin. For $f^n$ to
    have strictly more fixed points at the origin than $f$, the term $(a^n - 1)$
    from $f^n(z) - z = (a^n - 1)z + b_n z^{r+1} + \cdots = z((a^n - 1) + b_n z^r
    + \cdots)$ must vanish so that a higher power of $z$ takes over. Thus, we
    must have $a^n = 1$.
\end{proof}

% --- SECTION ------------------------------------------------------------------
% --- CRITICAL POINTS
% ----------------------------------------------------------
\section{Critical Points}
\label{sec:critical_points}

This section is devoted to the notion of critical points and critical values
even though their full significance for iteration will not be apparent until
Chapter~9 of \cite{Beardon1991}.

\begin{definition}[Critical point of a rational map]
    Given a rational map $R$. A point $z$ in $\ExtC$ is a \emph{critical point}
    of $R$ if $R$ fails to be injective in any neighborhood of $z$.
\end{definition}

\begin{remark}[Critical points in terms of valency]
    Observe that if $R$ is not constant, critical points are precisely the
    points at which $v_R(z) > 1$.
\end{remark}

\begin{definition}[Critical value of a rational map]
    Given a rational map $R$. A point $w$ in $\ExtC$ is a \emph{critical value}
    of $R$ if it is the image of some critical point; that is, if $w = R(z)$ for
    some critical point $z$ of $R$.
\end{definition}

\begin{remark}[Bijective restrictions of $R^{-1}$ for non-critical values]
    Given a rational function $R$ of degree $d \geq 1$, and $w$ a non-critical
    value of $R$. Then, $R^{-1}\{w\}$ consists of precisely $d$ distinct points,
    say $z_1, z_2, \dots, z_d$. As none of the $z_j$ are critical points, there
    are neighborhoods of $N$ of $w$, and $N_1, \ldots, N_d$ of $z_1, \ldots,
    z_d$ respectively, with $R$ acting as a bijection from each $N_j$ onto $N$
    (Figure~\ref{fig:bijective_restrictions_of_R_for_non_critical_value}). It
    follows that for each $j$, the restriction $R_j$ of $R$ to $N_j$ ($R: N_j
    \to N$) has an inverse
    \begin{equation*}
        R_j^{-1} : N \to N_j.
    \end{equation*}
\end{remark}

\begin{figure}
    \centering
    \resizebox{0.6\textwidth}{!}{\begin{tikzpicture}[
    >-Stealth,
    point/.style={circle, fill=black, inner sep=1pt},
    mapping arrow/.style={->, thick, shorten >=15pt, shorten <=15pt}
]

    % --- RIGHT SIDE: REGION N -------------------------------------------------
    % Center reference coordinate for the target region
    \coordinate (N_center) at (5, 0);
    
    % Draw Blob
    \pic [thick, scale=0.25, rotate=180] at (N_center) {blob7};

    \node[point, label=right:{$w$}] (w) at (N_center) {};
    \node[above right=0.85cm and 0.6cm of N_center] {$\mathcal{N}$};


    % --- LEFT SIDE: REGIONS N_1, N_2, N_d -------------------------------------
    
    % --- Region N_1 (Top Left) ------------------------------------------------
    \coordinate (N1_center) at (0, 2);
    
    % Draw Blob
    \pic [thick, scale=0.21, rotate=-5] at (N1_center) {blob1};

    \node[point, label=right:{$z_1$}] (z1) at (N1_center) {};
    \node[left=1.4cm of N1_center] {$\mathcal{N}_1$};

    % --- Region N_2 (Middle Left) ---------------------------------------------
    \coordinate (N2_center) at (0, 0);
    
    % Draw Blob
    \pic [thick, scale=0.2, rotate=-135] at (N2_center) {blob5};

    \node[point, label=right:{$z_2$}] (z2) at (N2_center) {};
    \node[left=1.4cm of N2_center] {$\mathcal{N}_2$};

    % Vertical Dots
    \node[font=\Large] at (0, -1.2) {$\vdots$};

    % --- Region N_d (Bottom Left) ---------------------------------------------
    \coordinate (Nd_center) at (0, -2.5);

    % Draw Blob
    \pic [thick, scale=0.2] at (Nd_center) {blob2};

    \node[point, label=right:{$z_d$}] (zd) at (Nd_center) {};
    \node[left=1.4cm of Nd_center] {$\mathcal{N}_d$};


    % --- MAPPING ARROWS -------------------------------------------------------
    
    % Arrow from N_1 to N
    \draw[mapping arrow] ($(N1_center) + (1.1, 0)$) -- ($(N_center) + (-0.5, 0.5)$) 
        node[midway, above=3pt] {$R_1$};
        
    % Arrow from N_2 to N
    \draw[mapping arrow] ($(N2_center) + (0.8, 0)$) -- ($(N_center) + (-0.8, -0.0)$) 
        node[midway, above=2pt] {$R_2$};
        
    % Arrow from N_d to N
    \draw[mapping arrow] ($(Nd_center) + (1.1, 0)$) -- ($(N_center) + (-0.5, -0.5)$) 
        node[midway, above=3pt] {$R_d$};

\end{tikzpicture}}
    \caption{Bijective restrictions of $R$ for a non-critical value $w$.}
    \label{fig:bijective_restrictions_of_R_for_non_critical_value}
\end{figure}

\begin{definition}[Branch of $R^{-1}$ at a non-critical value $w$]
    Given a rational function $R$ of degree $d \geq 1$, and $w$ a non-critical
    value of $R$. Then, we call each of the inverse maps $R_j^{-1} : N \to N_j$
    a \emph{branch} of $R^{-1}$ at $w$.
\end{definition}

By applying the Inverse Function Theorem
(Theorem~\ref{thm:inverse_function_theorem_in_C}), we know that $R$ is injective
on some neighborhood of any point in $\C$ at which $R'$ has neither a zero nor a
pole: thus for all but a finite set of $z$, $v_R(z) = 1$ (recall that locally
injective is the same as having valency 1
(Remark~\ref{rmk:local_injectivity_in_terms_of_valency})) and, consequently,
\begin{equation*}
    \sum_{z \in \ExtC} [v_R(z) - 1] < \infty.
\end{equation*}
This sum gives us a measure of the number of multiple roots of $R$ (and the
difficulties in defining $R^{-1}$), and its actual value is given by the
\emph{Riemann-Hurwitz relation}.

\begin{theorem}[Riemann-Hurwitz relation]
    \label{thm:Riemann_Hurwitz_relation}
    Given a non-constant rational map $R$. Then,
    \begin{equation*}
        \sum_{z \in \ExtC} [v_R(z) - 1] = 2\deg(R) - 2.
    \end{equation*}
\end{theorem}
\begin{proof}
    See \cite[44]{Beardon1991}.
\end{proof}

The general term in the sum is positive only when $z$ is a critical point of $R$
so, as the valency is an integer, this provides us with an estimate of the
number of critical points of $R$. For a non-constant polynomial $P$, we have
$v_P(\infty) = \deg(P)$, so, in the same way, we obtain an estimate of the
number of critical points of $P$ in $\C$. These two estimates are given in

\begin{corollary}[Bound on amount of critical points of a rational]
    A rational map of positive degree $d$ has at most $2d - 2$ critical points
    in $\ExtC$. A polynomial of positive degree $d$ has at most $d - 1$ critical
    points in $\C$.
\end{corollary}
\begin{proof}
    Given a non-constant rational map $R$ of degree $d$, for a point to be
    critical we must have $v_R(z) > 1$, so $v_R(z) - 1 \geq 1$. Thus, the number
    of critical points of $R$ in $\ExtC$ is at most $2d - 2$.

    In the case of polynomials, we have $v_P(\infty) = \deg(P) = d$, so the
    contribution of $\infty$ to the sum in the Riemann-Hurwitz relation is $d -
    1$. Thus, the number of critical points of $P$ in $\C$ is at most $(2d - 2)
    - (d - 1) = d - 1$.
\end{proof}

We shall see later that these estimates play a vital role in iteration theory.

\begin{definition}[Multiplicity of a critical point]
    Usually, we allocate a multiplicity to a critical point $z$: this
    multiplicity is $v_R(z) - 1$ and with this convention, a rational map of
    degree $d$ has exactly $2d - 2$ critical points.
\end{definition}

A topological proof of Theorem~\ref{thm:Riemann_Hurwitz_relation}, using the
Euler characteristic is available and has the advantage that, suitably modified,
it is valid for any analytic map between compact Riemann surfaces: this is given
in Chapter~\ref{chp:the_structure_of_the_Fatou_set}. However, at this point some
readers may prefer the previous more elementary proof.

We end this section with a simple result which locates the critical values of
the iterate $R^n$ in terms of the critical points of $R$.

\begin{theorem}[Critical values of $R^n$]
    Given a rational map $R$ and its set of critical points $C$. Then, the set
    of critical values of $R^n$ is
    \begin{equation*}
        R(C) \cup \cdots \cup R^n(C).
    \end{equation*}
\end{theorem}
\begin{proof}
    See \cite[45]{Beardon1991}.
\end{proof}

% --- SECTION ------------------------------------------------------------------
% --- A TOPOLOGY ON THE RATIONAL FUNCTIONS -------------------------------------
\section{A Topology on the Rational Functions}
\label{sec:a_topology_on_the_rational_functions}

To study the global dynamics and stability of families of rational maps, we must
equip the set of rational functions with a robust topological structure.

\begin{definition}[Spaces $\mathcal{C}$ and $\mathcal{R}$]
    Let $\mathcal{C}$ denote the class of all continuous maps from the extended
    complex plane into itself:
    \begin{equation*}
        \mathcal{C} = \{ f : \ExtC \longrightarrow \ExtC \mid f \text{ is continuous} \}.
    \end{equation*}
    We denote by $\mathcal{R}$ the specific subclass consisting of all
    non-constant and constant rational functions.
\end{definition}

\begin{definition}[Uniform Metric $\rho$]
    The \emph{metric of uniform convergence} $\rho : \mathcal{C} \times
    \mathcal{C} \longrightarrow \R^+ \cup \{0\}$ on the complex sphere is
    defined by:
    \begin{equation*}
        \rho(f, g) = \sup_{z \in \ExtC} \sigma_0(f(z), g(z)),
    \end{equation*}
    where $\sigma_0$ is the spherical metric.
\end{definition}

Since the chordal metric $\sigma$ and the spherical metric $\sigma_0$ only
differ by a bounded factor, we could equally well use the chordal metric here
(the metric would change, but the topology would not).

Now $\mathcal{R}$ is a closed subset of $\mathcal{C}$ because if the rational
functions $R_n$ converge uniformly to $R$ on the complex sphere, then $R$ is
analytic on the sphere and so it too is rational. More generally (and again we
could use either the chordal metric or the spherical metric), we have

\begin{theorem}[Uniform Limits Preserve Analyticity]
    Suppose that each function in a sequence $\{f_n\}$ is analytic in a domain
    $D \subseteq \ExtC$, and that $f_n$ converges uniformly on $D$ to $f$ with
    respect to the spherical metric $\sigma_0$. Then $f$ is also analytic in
    $D$.
\end{theorem}
\begin{proof}
    See \cite[46]{Beardon1991}.
\end{proof}

\begin{corollary}[$\mathcal{R}$ is Closed]
    The space of rational functions $\mathcal{R}$ is a closed subset of the
    continuous function space $\mathcal{C}$ under the topology induced by the
    uniform metric $\rho$.
\end{corollary}

With the topology established, we now investigate the topological behavior of
the degree function. Crucially, the degree of a rational map cannot change
abruptly under small uniform perturbations.

\begin{theorem}[Continuity of the Degree Function]
    \label{thm:continuity_of_degree_function}
    The degree map $\deg : \mathcal{R} \longrightarrow \{0, 1, 2, \dots\}$ is
    continuous. In particular, if a sequence of rational functions $R_n$
    converges uniformly on the Riemann sphere to a function $R$, then $R$ is
    rational and $\deg(R_n) = \deg(R)$ for all sufficiently large $n$.
\end{theorem}
\begin{proof}
    See \cite[46]{Beardon1991}.
\end{proof}

\begin{definition}[Space of Degree $n$ Maps]
    We denote by $\mathcal{R}_n$ the collection of all rational maps of degree
    exactly equal to $n$:
    \begin{equation*}
        \mathcal{R}_n = \{ R \in \mathcal{R} \mid \deg(R) = n \}.
    \end{equation*}
\end{definition}

As a consequence of this theorem
(Theorem~\ref{thm:continuity_of_degree_function}), the class $\mathcal{R}_n$ of
rational maps of degree $n$ is an open subset of $\mathcal{R}$ for it is the
inverse image of the open subset $\{n\}$ of $\Z$ under the continuous map
$\deg$. It is easy to see that each $\mathcal{R}_n$ is connected (for if $R$ and
$S$ lie in $\mathcal{R}_n$, we can simply \enquote{slide} the zeros and poles of
$R$ to those of $S$ while maintaining the same degree) so we obtain

\begin{corollary}[Components of $\mathcal{R}$]
    The subspaces $\mathcal{R}_n$ are the components of $\mathcal{R}$.
\end{corollary}

There are a few remarks that I've not included in the notes, because Beardon
points out that they can be safely ignored, in \cite[47]{Beardon1991}.